\documentclass[10pt,a4paper]{amsart}
\usepackage[margin=19mm]{geometry}
\usepackage{amsmath,amssymb,mathtools}
\usepackage[scr=rsfs,cal=euler]{mathalfa}
\usepackage{xfrac}
\usepackage[unicode,hidelinks,bookmarks=false]{hyperref}
\newcommand{\ie}{i.e.\ }
\newcommand{\cf}{cf.\ }
\newcommand{\resp}{respectively\ }
\newcommand{\Subsection}{\subsection}
\providecommand{\nobreakdash}{\nobreak\hbox{-}}
\setlength{\emergencystretch}{2em}
\multlinegap0pt
\usepackage{amssymb}
\usepackage{float}
\usepackage{xcolor}
\newtheorem{prop}{Proposition}
\title{The Test and Find Model}
\author{Michael Grabchak, Xingjie Li,, Isaac M. Sonin}
\date{Source: arXiv:2608.29390v1 (CC BY 4.0)}
\begin{document}
\maketitle

\noindent\emph{Source and licence.} The Test and Find Model. Version arXiv:2608.29390v1 (CC BY 4.0), distributed by arXiv under the Creative Commons Attribution 4.0 International licence, http://creativecommons.org/licenses/by/4.0/, as recorded in the arXiv licence metadata for this exact version and shown on the abstract page https://arxiv.org/abs/2608.29390v1. Full licence notice: \texttt{licenses/arxiv-2608.29390-CC-BY-4.0.txt}.\par\smallskip

\noindent\textit{Editorial scope.} Counted unit: the article's prop Prop1 (source0.tex lines 166--176). The article's complete original proof of it is delivered with it (source0.tex lines 178--192, one \texttt{proof} environment of 15 lines). No mathematics is written, repaired or re-derived: the statement and the proof are the article's own lines, in the article's own notation, and no retained line is substituted or re-flowed.  No further source-local context is needed: every cross-reference of every retained line resolves inside the two delivered spans.  Nothing was omitted from any retained span, so the package carries no omission record.  No retained line contains a citation, so the wrapper carries no bibliography and no bibliography entry.  No retained line contains a figure environment, an image-inclusion command or drawing code, so no image is delivered and the package declares no asset dependency.  Editorial formatting only, in addition: an AMS \texttt{amsart} preamble that restates the article's own declarations and macros used by the retained lines, namely the environments \texttt{prop} on separate counters, together with the standard packages those lines need.  The retained environments print in this document as Proposition 1 in order of appearance, whereas the article prints the same items with the numbering of its own full document.  The complete native source of the article as distributed in the arXiv e-print for this exact version is bundled unchanged as \texttt{sources/208847/source0.tex}.
\begin{prop}\label{Prop1} If $\gamma=(\gamma_1,\dots,\gamma_n) \in \Gamma$ and $s=(s_1,\dots,s_n)\in\Sigma$, then for any $\pi$
\begin{eqnarray}
p(s|\gamma)&=&\prod_{i: \gamma_i=1}a_i^{s_i}(1-a_i)^{1-s_i}\prod_{i: \gamma_i=0}b_i^{1-s_i}(1-b_i)^{s_i}, \label{eq: psg}\\
p(\gamma,s|\pi)&=&\pi(\gamma)p(s|\gamma),       \label{psgg}\\
p(s|\pi)&=&\sum_{\gamma\in \Gamma}\pi(\gamma)p(s|\gamma),       \label{pspi}\\
\theta(\gamma|s,\pi) &=&\frac{\pi(\gamma)p(s|\gamma)}{p(s|\pi)} , \label{pgs}\\
\alpha_i(s,\pi)&=& \sum_{\gamma\in\Gamma: \gamma_i=0}\theta(\gamma |s,\pi)=
\frac{\sum\limits_{\gamma\in\Gamma: \gamma_i=0} \pi(\gamma)p(s|\gamma)}{\sum_{\gamma\in \Gamma}\pi(\gamma)p(s|\gamma)}, \mbox{ and }\label{eq: alpha i}\\
\sum_{i=1}^n\alpha_i(s,\pi) &=& n-k. \label{snk}
\end{eqnarray}
\end{prop}

\begin{proof}
Since the $S_i$'s are conditionally independent given $T=\gamma$,
$$
p(s|\gamma)=\prod_{i=1}^{n}p(s_i|\gamma)=\prod_{i: \gamma_i=1}p(s_i|\gamma)\prod_{i: \gamma_i=0}p(s_i|\gamma).
$$
Furthermore, given $T=\gamma$, the random variable $S_i$ has a Bernoulli distribution with pmf $p(x)=p^x(1-p)^{1-x}$ for $x=0,1$, where $p=a_i$ if $\gamma_i=1$ and $p=1-b_i$ if $\gamma_i=0$. From here \eqref{eq: psg} follows directly. Next, \eqref{psgg} follows from the definition of conditional probability, \eqref{pspi} from the law of total probability, \eqref{pgs} from Bayes' theorem, and \eqref{eq: alpha i} from the fact that $\alpha_i(s,\pi)= P(T_i=0 |S=s,\pi)=\sum_{\gamma\in\Gamma: \gamma_i=0}P(T=\gamma|S=s,\pi)$ combined with \eqref{pgs} and \eqref{pspi}. Finally, to show \eqref{snk}, let $1_{[T_i=0]}$ denote the indicator function on event $[T_i=0]$ and note that $\sum_{i=1}^n 1_{[T_i=0]}$ is the number of empty boxes, which is $n-k$. From the definition of $\alpha_i(s,\pi)$ to have
\[
\begin{aligned}
\sum_{i=1}^n\alpha_i(s,\pi)=&
\sum_{i=1}^n P(T_i=0 |S=s,\pi)\\
=&\mathrm E\left[\sum_{i=1}^n 1_{[T_i=0]}\Big{|}S= s,\pi\right] = \mathrm E\left[ n-k\Big{|} S=s,\pi\right] = n-k,
\end{aligned}
\]
which completes the proof.
\end{proof}

\end{document}
